% Generated by scripts/assemble.py; corrections belong in source/corrections.json.
\examhead{67}{太原理工大学 2025 年数学分析试卷}{67}
% p068-b06; source PDF page 68
\blockanchor{p068-b06}
\question{1}求极限 $\lim_{x\to 0}\frac{(1 + \sin^2x)^{1011} - (\cos x)^{2026}}{\sqrt{1 + 2x^2} - 1}.$\par
% p068-b07; source PDF page 68
\blockanchor{p068-b07}
\question{2}求极限  $\lim_{n\to\infty}\left[(n+2025)\left(\frac{1}{n^{2}+1}+\frac{1}{n^{2}+2^{2}}+\cdots+\frac{1}{n^{2}+n^{2}}\right)\right]$ .\par
% p068-b08; source PDF page 68
\blockanchor{p068-b08}
\question{3}求不定积分 $\int \frac{1 + \sin x}{\sin x(1 + \cos x)} \mathrm{d}x.$\par
% p068-b09; source PDF page 68
\blockanchor{p068-b09}
\question{4}求累次积分 $\int_0^1 x^2\mathrm{d}x\int_x^1 e^{-y^2}\mathrm{d}y.$\par
% p068-b10; source PDF page 68
\par\Needspace{6\baselineskip}
\blockanchor{p068-b10}
\question{5}计算曲面积分\par
\nopagebreak[4]
% p068-b11; source PDF page 68
\blockanchor{p068-b11}
\displayformula{I = \iint_ {\Sigma} 4 x z \mathrm{d} y \mathrm{d} z - 2 y z \mathrm{d} z \mathrm{d} x + (1 - z ^ {2}) \mathrm{d} x \mathrm{d} y.}
% p068-b12; source PDF page 68
\blockanchor{p068-b12}
其中 $S$ 是旋转曲面 $z = e^{\sqrt{x^2 + y^2}}$ 介于 $z = 1$ 和 $z = e^2$ 的部分, 取下侧.\par
% p068-b13; source PDF page 68
\blockanchor{p068-b13}
\question{6}求幂级数 $\sum_{n=1}^{\infty} (-1)^{n-1} n^2 x^n$ 的收敛域与和函数.\par
% p068-b14; source PDF page 68
\blockanchor{p068-b14}
\question{7}求 $I(a) = \int_{0}^{+\infty} e^{-x^2} \cos(ax) \, \mathrm{d}x$.\par
% p068-b15; source PDF page 68
\par\Needspace{6\baselineskip}
\blockanchor{p068-b15}
\question{8}讨论函数\par
\nopagebreak[4]
% p068-b16; source PDF page 68
\blockanchor{p068-b16}
\displayformula{f (x, y) = \left\{ \begin{array}{l l} \frac {x ^ {3} - y ^ {3}}{x ^ {2} + y ^ {2}}, & x ^ {2} + y ^ {2} \neq 0; \\ 0, & x ^ {2} + y ^ {2} = 0. \end{array} \right.}
% p068-b17; source PDF page 68
\blockanchor{p068-b17}
在原点 $(0,0)$ 处的连续性, 偏导数的存在性, 可微性.\par
% p068-b18; source PDF page 68
\par\Needspace{6\baselineskip}
\blockanchor{p068-b18}
\question{9}$f(x)$ 在 $[0, +\infty)$ 上非负连续，且\par
\nopagebreak[4]
% p068-b19; source PDF page 68
\blockanchor{p068-b19}
\displayformula{f (x) \leq c + \int_ {0} ^ {x} f (t) \mathrm{d} t, x \geq 0.}
% p068-b20; source PDF page 68
\blockanchor{p068-b20}
其中 c > 0 为常数. 证明:  $f(x) \leq ce^{x}, x \geq 0$ .\par
% p068-b21; source PDF page 68
\par\Needspace{6\baselineskip}
\blockanchor{p068-b21}
\question{10}证明:\par
\nopagebreak[4]
% p068-b22; source PDF page 68
\blockanchor{p068-b22}
\displayformula{\frac {x}{1 + x} <   \ln (1 + x) <   x, x > 0. \tag {1}}
% p068-b23; source PDF page 68
\blockanchor{p068-b23}
(2) $\sum_{n=1}^{\infty} \frac{(-1)^n}{n} \cdot \frac{\ln(1 + x^n)}{x^n}$ 在 $(0, +\infty)$ 上一致收敛.\par
% p068-b24; source PDF page 68
\blockanchor{p068-b24}
\question{11}$f(x)$ 在 $[0,1]$ 上连续，$(0,1)$ 上可导，且 $f(0) = f(1) = 1$。证明：\par
% p069-b00; source PDF page 69
\blockanchor{p069-b00}
(1) 存在 $\eta \in (0,1)$, 使得 $e^{\eta}[f(\eta) + f'(\eta)] = e - 1$.\par
% p069-b01; source PDF page 69
\blockanchor{p069-b01}
(2) 存在 $\xi \in (0,1)$, 使得 $e^{\eta - \xi}[f(\eta) + f'(\eta)] = 1$.\par
% p069-b02; source PDF page 69
\blockanchor{p069-b02}
\question{12}设 $\{f_n(x)\}$ 在 $[a, b]$ 上一致收敛于 $f(x)$, 且对每一个正整数 $n$, $f_n(x)$ 均在 $[a, b]$ 上连续. 证明:\par
% p069-b03; source PDF page 69
\blockanchor{p069-b03}
(1) $f(x)$ 在 $[a, b]$ 上一致连续.\par
% p069-b04; source PDF page 69
\par\Needspace{6\baselineskip}
\blockanchor{p069-b04}
(2) 对任意的 $\varepsilon > 0$, 存在 $\delta > 0$, 对任意的 $x', x'' \in [a, b]$, 只要 $|x' - x''| < \delta$, 对任意的正整数 $n$, 都有\par
\nopagebreak[4]
% p069-b05; source PDF page 69
\blockanchor{p069-b05}
\displayformula{\left| f _ {n} (x ^ {\prime}) - f _ {n} (x ^ {\prime \prime}) \right| <   \varepsilon .}
